% GENERATED FILE. Edit canonical lessons, then run npm run generate:books.
% canonical-source-hash: fnv1a64:be3008de8645919a
% canonical-lessons: HI-C523-caku, HI-C523-chiilna, HI-C523-kutna, HI-C523-burakna, HI-C523-dalnagrind

\chapter{Peel, Pound, Strain}
\label{ch:hi-peel-pound-strain}
\hlchaptermodality{523}

\begin{chapteropening}
\textbf{By the end of this chapter:} \emph{I can describe getting food ready to cook.}

Complete the last lesson of chapter 523: I can describe getting food ready to cook.
\end{chapteropening}

\section[\texorpdfstring{\hi{चाकू} (cākū)}{चाकू (cākū)}]{\hi{चाकू} (cākū) — a knife, a pocket knife}

\label{lesson:HI-C523-caku}



\begin{quote}
Before the new one: say the Hindi for to seethe, then the Hindi for to drain off (liquid).
\end{quote}

\subsection*{You'll want to know: \hi{चाकू}}
\textbf{\hi{चाकू}} — \emph{cākū} — \textquotedblleft{}a knife, a pocket knife\textquotedblright{}.

\begin{quote}
\hi{सब्ज़ी} \hi{काटने} \hi{के} \hi{लिए} \hi{एक} \hi{तेज़} \hi{चाकू} \hi{दीजिए।}
\end{quote}

\begin{quote}
\emph{Sabzī kāṭne ke liye ek tez cākū dījiye.}
\end{quote}

\begin{quote}
Please give me a sharp knife to cut the vegetables.
\end{quote}

Masculine.

\begin{cousinweb}
\textbf{\hi{चाकू}} (\emph{cākū}) is usually traced to Persian \emph{chāqū}, \textquotedblleft{}a knife\textquotedblright{}.
\end{cousinweb}

\begin{grammarlens}[title={What you've built}]
One more word for describing what you lived and felt.
\end{grammarlens}

\subsection*{Your turn}
\begin{itemize}
\raggedright
  \item \emph{Say it:} \emph{cākū}
  \item \emph{Say it:} \emph{cākū}, once more
  \item \emph{Say it:} \emph{Sabzī kāṭne ke liye ek tez cākū dījiye.} — Please give me a sharp knife to cut the vegetables.
  \item \emph{Recall:} say the Hindi for to seethe, then the Hindi for to drain off (liquid), then say \emph{cākū} again
\end{itemize}

\begin{tcolorbox}[breakable,colback=teal!4,colframe=teal!35!black,title={Before you move on}]
What does \hi{चाकू} mean? (\textquotedblleft{}A knife, a pocket knife\textquotedblright{}.) Say it once more.
\end{tcolorbox}
\section[\texorpdfstring{\hi{छीलना} (chīlnā)}{छीलना (chīlnā)}]{\hi{छीलना} (chīlnā) — to peel}

\label{lesson:HI-C523-chiilna}



\begin{quote}
Before the new one: say the Hindi for to drain off (liquid), then the Hindi for a pocket knife.
\end{quote}

\subsection*{You'll want to know: \hi{छीलना}}
\textbf{\hi{छीलना}} — \emph{chīlnā} — \textquotedblleft{}to peel\textquotedblright{}.

\begin{quote}
\hi{मीरा} \hi{आलू} \hi{छील} \hi{रही} \hi{है।}
\end{quote}

\begin{quote}
\emph{Mīrā ālū chīl rahī hai.}
\end{quote}

\begin{quote}
Mira is peeling potatoes.
\end{quote}

Past, with \hi{ने} (he or she …): \hi{उसने} \hi{केला} \hi{छीला}, \emph{usne kelā chīlā}.

\begin{cousinweb}
\textbf{\hi{छीलना}} (\emph{chīlnā}) is a tadbhava, inherited from \emph{*chilla}, \textquotedblleft{}torn, cut\textquotedblright{}, a form CDIAL reconstructs from its descendants (CDIAL 5051). CDIAL gives the Hindi verb as \textquotedblleft{}to peel\textquotedblright{}.
\end{cousinweb}

\begin{grammarlens}[title={What you've built}]
One more word for describing what you lived and felt.
\end{grammarlens}

\subsection*{Your turn}
\begin{itemize}
\raggedright
  \item \emph{Say it:} \emph{chīlnā}
  \item \emph{Say it:} \emph{chīlnā}, once more
  \item \emph{Say it:} \emph{Mīrā ālū chīl rahī hai.} — Mira is peeling potatoes.
  \item \emph{Recall:} say the Hindi for to drain off (liquid), then the Hindi for a pocket knife, then say \emph{chīlnā} again
\end{itemize}

\begin{tcolorbox}[breakable,colback=teal!4,colframe=teal!35!black,title={Before you move on}]
What does \hi{छीलना} mean? (\textquotedblleft{}To peel\textquotedblright{}.) Say it once more.
\end{tcolorbox}
\section[\texorpdfstring{\hi{कूटना} (kūṭnā)}{कूटना (kūṭnā)}]{\hi{कूटना} (kūṭnā) — to pound}

\label{lesson:HI-C523-kutna}



\begin{quote}
Before the new one: say the Hindi for a pocket knife, then the Hindi for to peel.
\end{quote}

\subsection*{You'll want to know: \hi{कूटना}}
\textbf{\hi{कूटना}} — \emph{kūṭnā} — \textquotedblleft{}to pound\textquotedblright{}.

\begin{quote}
\hi{उसने} \hi{अदरक} \hi{को} \hi{पत्थर} \hi{से} \hi{कूटा।}
\end{quote}

\begin{quote}
\emph{Usne adrak ko patthar se kūṭā.}
\end{quote}

\begin{quote}
She pounded the ginger with a stone.
\end{quote}

Past, with \hi{ने} (he or she …): \hi{उसने} \hi{लहसुन} \hi{कूटा}, \emph{usne lahsun kūṭā}.

\begin{cousinweb}
\textbf{\hi{कूटना}} (\emph{kūṭnā}) is a tadbhava, inherited from Sanskrit \emph{kuṭṭayati}, \textquotedblleft{}pounds\textquotedblright{} (CDIAL 3241).
\end{cousinweb}

\begin{grammarlens}[title={What you've built}]
One more word for describing what you lived and felt.
\end{grammarlens}

\subsection*{Your turn}
\begin{itemize}
\raggedright
  \item \emph{Say it:} \emph{kūṭnā}
  \item \emph{Say it:} \emph{kūṭnā}, once more
  \item \emph{Say it:} \emph{Usne adrak ko patthar se kūṭā.} — She pounded the ginger with a stone.
  \item \emph{Recall:} say the Hindi for a pocket knife, then the Hindi for to peel, then say \emph{kūṭnā} again
\end{itemize}

\begin{tcolorbox}[breakable,colback=teal!4,colframe=teal!35!black,title={Before you move on}]
What does \hi{कूटना} mean? (\textquotedblleft{}To pound\textquotedblright{}.) Say it once more.
\end{tcolorbox}
\section[\texorpdfstring{\hi{बुरकना} (buraknā)}{बुरकना (buraknā)}]{\hi{बुरकना} (buraknā) — to sprinkle (a powder)}

\label{lesson:HI-C523-burakna}



\begin{quote}
Before the new one: say the Hindi for to peel, then the Hindi for to pound.
\end{quote}

\subsection*{You'll want to know: \hi{बुरकना}}
\textbf{\hi{बुरकना}} — \emph{buraknā} — \textquotedblleft{}to sprinkle (a powder)\textquotedblright{}.

\begin{quote}
\hi{माँ} \hi{ने} \hi{सब्ज़ी} \hi{पर} \hi{थोड़ा} \hi{नमक} \hi{बुरका।}
\end{quote}

\begin{quote}
\emph{Māṁ ne sabzī par thoṛā namak burkā.}
\end{quote}

\begin{quote}
Mother sprinkled a little salt over the vegetables.
\end{quote}

Past, with \hi{ने} (he or she …): \hi{उसने} \hi{हलवे} \hi{पर} \hi{चीनी} \hi{बुरकी}, \emph{usne halve par cīnī burkī}.

\begin{cousinweb}
\textbf{\hi{बुरकना}} (\emph{buraknā}) is a tadbhava, inherited from \emph{*būra}, \textquotedblleft{}powder\textquotedblright{}, a form CDIAL reconstructs from its descendants (CDIAL 9298). CDIAL gives the Hindi verb as \textquotedblleft{}to sprinkle\textquotedblright{}.
\end{cousinweb}

\begin{grammarlens}[title={What you've built}]
One more word for describing what you lived and felt.
\end{grammarlens}

\subsection*{Your turn}
\begin{itemize}
\raggedright
  \item \emph{Say it:} \emph{buraknā}
  \item \emph{Say it:} \emph{buraknā}, once more
  \item \emph{Say it:} \emph{Māṁ ne sabzī par thoṛā namak burkā.} — Mother sprinkled a little salt over the vegetables.
  \item \emph{Recall:} say the Hindi for to peel, then the Hindi for to pound, then say \emph{buraknā} again
\end{itemize}

\begin{tcolorbox}[breakable,colback=teal!4,colframe=teal!35!black,title={Before you move on}]
What does \hi{बुरकना} mean? (\textquotedblleft{}To sprinkle (a powder)\textquotedblright{}.) Say it once more.
\end{tcolorbox}
\section[\texorpdfstring{\hi{दलना} (dalnā)}{दलना (dalnā)}]{\hi{दलना} (dalnā) — to grind coarsely; to split (pulses)}

\label{lesson:HI-C523-dalnagrind}



\begin{quote}
Before the new one: say the Hindi for to pound, then the Hindi for to sprinkle (a powder).
\end{quote}

\subsection*{You'll want to know: \hi{दलना}}
\textbf{\hi{दलना}} — \emph{dalnā} — \textquotedblleft{}to grind coarsely; to split (pulses)\textquotedblright{}.

\begin{quote}
\hi{माँ} \hi{ने} \hi{घर} \hi{पर} \hi{दाल} \hi{दली।}
\end{quote}

\begin{quote}
\emph{Māṁ ne ghar par dāl dalī.}
\end{quote}

\begin{quote}
Mother split the pulses at home.
\end{quote}

Past, with \hi{ने} (he or she …): \hi{उसने} \hi{चावल} \hi{दले}, \emph{usne cāval dale}.

\begin{cousinweb}
\textbf{\hi{दलना}} (\emph{dalnā}) is a tadbhava, inherited from Sanskrit \emph{dalati}, \textquotedblleft{}bursts, splits\textquotedblright{} (CDIAL 6216). CDIAL gives the Hindi verb as \textquotedblleft{}to grind coarsely\textquotedblright{}.
\end{cousinweb}

\begin{grammarlens}[title={What you've built}]
One more word for describing what you lived and felt.
\end{grammarlens}

\subsection*{Your turn}
\begin{itemize}
\raggedright
  \item \emph{Say it:} \emph{dalnā}
  \item \emph{Say it:} \emph{dalnā}, once more
  \item \emph{Say it:} \emph{Māṁ ne ghar par dāl dalī.} — Mother split the pulses at home.
  \item \emph{Recall:} say the Hindi for to pound, then the Hindi for to sprinkle (a powder), then say \emph{dalnā} again
\end{itemize}

\begin{tcolorbox}[breakable,colback=teal!4,colframe=teal!35!black,title={Before you move on}]
What does \hi{दलना} mean? (\textquotedblleft{}To grind coarsely; to split (pulses)\textquotedblright{}.) Say it once more.
\end{tcolorbox}
